% 引力波天文学（综述）
% license CCBYSA3
% type Wiki

本文根据 CC-BY-SA 协议转载翻译自维基百科\href{https://en.wikipedia.org/wiki/Gravitational-wave_astronomy}{相关文章}。


引力波天文学是天文学的一个分支，主要研究由天体物理源发出的引力波的探测与分析。\(^\text{[1]}\)

引力波是时空中极其微小的扭曲或涟漪，由大质量天体的加速运动所引起。它们产生于灾变性事件之中，例如双黑洞并合、双中子星合并、超新星爆炸，以及大爆炸后早期宇宙中的物理过程。研究引力波为观测宇宙提供了一种全新的方式，使人们能够深入了解物质在极端条件下的行为。类似于电磁辐射（如光波、无线电波、红外辐射和 X 射线）通过电磁场涨落传播能量的方式，引力辐射则是由相对较弱的引力场涨落所产生的能量传播。引力波的存在最早由奥利弗·赫维赛德（Oliver~Heaviside）于 1893 年提出，随后在 1905 年由昂利·庞加莱（Henri~Poincaré）作为电磁波的引力学对应物加以推测，最终在 1916 年被阿尔伯特·爱因斯坦（Albert~Einstein）根据广义相对论所预言。

1978 年，罗素·阿伦·赫尔斯（Russell~Alan~Hulse）与约瑟夫·胡顿·泰勒（Joseph~Hooton~Taylor~Jr.）通过观测两颗互相绕转的中子星，首次提供了引力波存在的实验性证据。二人因此工作获得了 1993 年诺贝尔物理学奖。2015 年，距爱因斯坦的预言近一个世纪后，人类首次直接探测到了由双黑洞并合产生的引力波信号，从而确认了这一长期难以捉摸的现象的存在，并开启了天文学的新纪元。此后，科学家又陆续探测到多起包括双黑洞并合、中子星碰撞及其他剧烈宇宙事件在内的引力波信号。如今，引力波的探测主要依靠激光干涉测量技术（laser interferometry），该技术通过测量两条相互垂直臂长随波动变化的微小差异来捕捉引力波。诸如 LIGO（Laser~Interferometer~Gravitational-wave~Observatory，激光干涉引力波天文台）、Virgo 与 KAGRA（Kamioka~Gravitational~Wave~Detector，神冈引力波探测器）等观测站都采用此项技术来探测来自遥远宇宙事件的微弱信号。LIGO 的联合创始人巴里·C·巴里什（Barry~C.~Barish）、基普·S·索恩（Kip~S.~Thorne）与赖纳·魏斯（Rainer~Weiss）因其在引力波天文学领域的开创性贡献而共同获得了 2017 年诺贝尔物理学奖。

当利用电磁波观测遥远的天体时，散射、吸收、反射、折射等多种物理过程会导致信息丢失。  
宇宙中存在若干光子只能部分穿透的区域，例如星云内部、银河系核心的致密尘埃云以及黑洞附近的区域等。与电磁天文学相互补充的引力波天文学，有潜力以更高的分辨率研究宇宙。在一种被称为“多信使天文学”（multi-messenger astronomy）的方法中，引力波数据与不同波段的电磁观测数据相结合，以获得对天体物理现象更全面的理解。引力波天文学不仅有助于揭示早期宇宙的状态、检验引力理论、探测暗物质与暗能量的分布，还能帮助确定哈勃常数（Hubble constant），即描述宇宙加速膨胀速率的重要参数。这些研究为超越标准模型（Beyond the Standard Model, BSM）的物理学打开了新的大门。


当前该领域仍面临多项挑战，包括噪声干扰、超高灵敏度探测设备的缺乏以及低频引力波信号的探测难题。地面探测器受到环境扰动引起的地震振动影响，同时其探测臂长度也受到地球曲率的限制。未来，引力波天文学将致力于研发升级型探测装置与下一代观测台，同时探索基于太空的探测方案，例如 LISA（Laser~Interferometer~Space~Antenna，激光干涉空间天线）。  
LISA 将能够探测来自银河系核心致密超大质量黑洞、原初黑洞等遥远源的信号，以及诸如双白矮星并合与早期宇宙来源的低频引力波信号。\(^\text{[2]}\)
\subsection{引言}
引力波是由轨道双体系统中加速运动的质量所产生的引力强度扰动波，它们以光速从源处向外传播。这种波动最早由奥利弗·赫维赛德（Oliver~Heaviside）于 1893 年提出，后又由昂利·庞加莱（Henri~Poincaré）于 1905 年进一步提出，认为它们类似于电磁波，但属于引力的对应形式。


1916 年，阿尔伯特·爱因斯坦（Albert~Einstein）在其广义相对论的框架下正式预言了引力波的存在，并将其描述为时空中的涟漪。然而，爱因斯坦后来一度拒绝接受引力波的存在。\(^\text{[3]}\)引力波以引力辐射（gravitational~radiation）的形式传输能量，这种辐射是一种类似于电磁辐射的辐射能。经典力学中的牛顿万有引力定律并不包含引力波的存在，因为该定律假设物理作用是瞬时传播（即传播速度无限大）的——这恰好揭示了牛顿力学在解释相对论相关现象方面的局限性。

引力波存在的第一个间接证据于 1974 年发现，来自对赫尔斯–泰勒双星脉冲星（Hulse–Taylor~binary~pulsar）轨道衰变的观测，其衰减速率与广义相对论预测的因引力辐射能量损失所导致的结果一致。1993 年，罗素·A·赫尔斯（Russell~A.~Hulse）与约瑟夫·胡顿·泰勒（Joseph~Hooton~Taylor~Jr.）因这一发现共同获得诺贝尔物理学奖。

直到 2015 年，人类才首次直接观测到引力波信号——该信号来自两颗黑洞并合事件，被美国路易斯安那州利文斯顿（Livingston,~Louisiana）与华盛顿州汉福德（Hanford,~Washington）的 LIGO（Laser~Interferometer~Gravitational-wave~Observatory，激光干涉引力波天文台）探测器接收到。  
2017 年诺贝尔物理学奖随后授予赖纳·魏斯（Rainer~Weiss）、基普·索恩（Kip~Thorne）与巴里·巴里什（Barry~Barish），以表彰他们在引力波直接探测中的奠基性贡献。

在引力波天文学（gravitational-wave~astronomy）中，科学家通过观测引力波来推断其源的物理性质。  
能够通过这种方式研究的源包括：由白矮星、中子星或黑洞组成的双星系统；超新星爆发事件；以及宇宙大爆炸后早期宇宙形成阶段的物理过程。
\subsubsection{仪器与挑战}
由于不同探测器在技术参数与灵敏度上的差异，探测器间的协作有助于收集互补且极具价值的数据。目前存在多台地面激光干涉仪，长度跨越数英里或数公里，包括：位于美国华盛顿州与路易斯安那州的两座 LIGO 探测器；位于意大利欧洲引力观测台（European~Gravitational~Observatory）的 Virgo 探测器；位于德国的 GEO600；  
以及位于日本的神冈引力波探测器 KAGRA（Kamioka~Gravitational~Wave~Detector）。迄今为止，LIGO、Virgo 与 KAGRA 已开展多次联合观测；而 GEO600 由于其设备灵敏度相对较低，目前主要用于实验性测试与技术验证，近年未参与其他观测台的联合运行。
\begin{figure}[ht]
\centering
\includegraphics[width=6cm]{./figures/e66390bc2706521f.png}
\caption{} \label{fig_Gras_1}
\end{figure}